\documentclass[11pt]{article}
\usepackage{preamble}
\setlength{\parskip}{4pt}

\begin{document}

\daybanner{AP Statistics \textperiodcentered\ Unit 2 \textperiodcentered\ Topic 2.9 \textperiodcentered\ B: 2026-11-02 \textperiodcentered\ E: 2026-11-23 \textperiodcentered\ TEACHER KEY}{Same Mean, Different Deadline Risk}

\sectionbanner{CED TETHER}

{\small
\begin{itemize}[leftmargin=1.5em,itemsep=2pt,topsep=2pt]
  \item \textbf{Skill 3.B} --- Determine parameters for probability distributions.
  \item \textbf{Skill 4.B} --- Interpret statistical calculations and findings to assign meaning or assess a claim.
  \item \textbf{EU VAR-5} --- Probability distributions may be used to model variation in populations.
  \item \textbf{LO VAR-5.C} --- Calculate parameters for a discrete random variable. [Skill 3.B]
  \item \textbf{EK VAR-5.C.1} --- A numerical value measuring a characteristic of a population or the distribution of a random variable is known as a parameter, which is a single, fixed value.
  \item \textbf{EK VAR-5.C.2} --- The mean, or expected value, for a discrete random variable X is $\mu_X$ = $\sum$ $x_i$ $\cdot$ P($x_i$).
  \item \textbf{EK VAR-5.C.3} --- The standard deviation for a discrete random variable X is $\sigma_X$ = $\surd$[$\sum$ ($x_i$ - $\mu_X$)\textsuperscript{2} $\cdot$ P($x_i$)].
  \item \textbf{LO VAR-5.D} --- Interpret parameters for a discrete random variable. [Skill 4.B]
  \item \textbf{EK VAR-5.D.1} --- Parameters for a discrete random variable should be interpreted using appropriate units and within the context of a specific population.
\end{itemize}
}
\textbf{Essential question:} When two probability models have the same mean, how can standard deviation change the decision?\par
\textbf{DOK-3 skill:} 4.B \textperiodcentered\ \textbf{FRQ pattern:} {\small\texttt{same-\allowbreak{}mean-\allowbreak{}different-\allowbreak{}sd-\allowbreak{}which-\allowbreak{}parameter-\allowbreak{}answers-\allowbreak{}the-\allowbreak{}question}}\par
\textbf{DOK rationale:} Part (c) weighs competing claims using the day's numerical or graphical evidence and explains a limitation or consequence; the judgment requires linked reasoning beyond a calculation.\par

\frameworkphaseheader{Do Now (first take) $\rightarrow$ Explore (video + follow-along) $\rightarrow$ Exit (finish + turn in)}{1 $\rightarrow$ 3}{45}{%
\begin{itemize}[leftmargin=1.5em,itemsep=2pt,topsep=2pt]
  \item Show the board at the bell and collect a one-sentence first take without confirming a claim.
  \item Run the named video follow-along, then allow about 10 minutes for parts (a)--(c).
  \item Ask for evidence linked to a claim. Collect the sheet and hand-score part (c) with E/P/I.
\end{itemize}
}{%
\begin{itemize}[leftmargin=1.5em,itemsep=2pt,topsep=2pt]
  \item Choose a route and cite one model feature.
  \item Complete the video follow-along about means, standard deviations, and their interpretations.
  \item Finish (a)--(c), revise the first take in the exit line, and turn in the sheet.
\end{itemize}
}{%
\begin{itemize}[leftmargin=1.5em,itemsep=2pt,topsep=2pt]
  \item Which number or visible feature supports your claim?
  \item Why does the other claim go beyond that evidence?
  \item What would you tell someone who chose the other claim?
\end{itemize}
}{Circulate and ask students to connect their choice to evidence. Score part (c) using the three required reasoning elements; do not require dropped extension work.}

\begin{callout}{\IconWarn\ WATCH FOR}{calloutred}
\begin{itemize}[leftmargin=1.5em,itemsep=2pt,topsep=2pt]
  \item A choice with copied numbers but no explanation of how they support it.
  \item Treating a model or average as a guaranteed individual outcome.
\end{itemize}
\end{callout}

\sectionbanner{SCORING PART (c) --- E / P / I}

\textbf{Expected elements}
\begin{itemize}[leftmargin=1.5em,itemsep=2pt,topsep=2pt]
  \item \textbf{reasoning-1} (required) --- Chooses A using standard deviations 3 versus 12 minutes
  \item \textbf{reasoning-2} (required) --- Connects 0/8 versus 2/8 above 40 to the customer's deadline
  \item \textbf{reasoning-3} (required) --- Explains that equal means do not imply equal variability or late risk
\end{itemize}

{\small\renewcommand{\arraystretch}{1.25}
\noindent\begin{tabular}{|p{0.5in}|p{5.9in}|}
\hline
\textbf{E} & Includes all three required reasoning elements above, with correct contextual evidence. \\ \hline
\textbf{P} & Includes at least one substantive correct reasoning link above, but omits or misstates another required element. \\ \hline
\textbf{I} & Includes no substantive correct reasoning link above; an unsupported choice or copied number alone is insufficient. \\ \hline
\end{tabular}}

\textbf{Common mistakes}
\begin{itemize}[leftmargin=1.5em,itemsep=2pt,topsep=2pt]
  \item Treating standard deviation as minutes added to every delivery
  \item Claiming every Route B time is exactly 12 minutes from its mean
  \item Using equal means alone to infer equal chances of missing the deadline
\end{itemize}

\clearpage
\focusbanner{TODAY'S PROBLEM --- annotated}

Suppose a dispatcher models delivery time using eight equally likely outcomes for each route, shown in the dotplots. An analyst reports that both routes have mean 30 minutes, with $\sigma_A=3$ minutes and $\sigma_B=12$ minutes.\par\smallskip The dispatcher says, \emph{The routes have the same mean, so it does not matter which one a customer gets.} A customer has 40 minutes before an appointment and places a high cost on being late.\par For the eight outcomes, the sums of squared distances from 30 are 72 for Route A and 1,152 for Route B.\par
\begin{center}
\begin{tikzpicture}
\begin{axis}[
  width=0.90\linewidth,
  height=1.0in,
  xmin=5,
  xmax=55,
  xtick={10,15,20,25,30,35,40,45,50},
  ymin=0,
  ymax=3,
  ytick=\empty,
  axis y line=none,
  axis x line=bottom,
  title={\textbf{Route A}},
  xlabel={Delivery time (minutes)},
  tick label style={font=\small},
  label style={font=\small},
  title style={font=\small}
]
\addplot[only marks,mark=*,mark size=2.4pt,color=dayblue] coordinates
  {(25,1) (27,1) (29,1) (29,2) (31,1) (31,2) (33,1) (35,1)};
\end{axis}
\end{tikzpicture}\par\smallskip
\begin{tikzpicture}
\begin{axis}[
  width=0.90\linewidth,
  height=1.0in,
  xmin=5,
  xmax=55,
  xtick={10,15,20,25,30,35,40,45,50},
  ymin=0,
  ymax=3,
  ytick=\empty,
  axis y line=none,
  axis x line=bottom,
  title={\textbf{Route B}},
  xlabel={Delivery time (minutes)},
  tick label style={font=\small},
  label style={font=\small},
  title style={font=\small}
]
\addplot[only marks,mark=*,mark size=2.4pt,color=dayblue] coordinates
  {(10,1) (18,1) (26,1) (26,2) (34,1) (34,2) (42,1) (50,1)};
\end{axis}
\end{tikzpicture}
\end{center}
\begin{firsttakebox}{FIRST TAKE (5 min)}
Which route would you assign to this customer? Cite one number or dotplot feature.\par
\answer{Not graded. Route B's very fast outcomes may attract some students before deadline risk is made explicit.}
\end{firsttakebox}

\textit{Rules callout ``MEAN AND STANDARD DEVIATION (keep on the page)'' is printed on the student sheet.}\par\medskip

\dokbadge{1}~\textbf{(a)} Count the outcomes longer than 40 minutes for each route.\par
\answer{Route A has no outcomes longer than 40 minutes. Route B has two: 42 and 50 minutes.}

\dokbadge{2}~\textbf{(b)} Use the given sums of squared distances to verify both population standard deviations. Interpret what one of them means in minutes.\par
\answer{$\sigma_A=\sqrt{72/8}=3$ minutes; $\sigma_B=\sqrt{1152/8}=12$ minutes. Route A's modeled delivery times typically differ from their 30-minute mean by about 3 minutes.}

\dokbadge{3}~\textbf{(c)} Which route fits this customer's priority? Use both standard deviations and the chance of exceeding 40 minutes. Explain why the dispatcher's mean-only comparison can mislead this customer. Write 3--4 sentences.\par
\answer{Choose Route A: its standard deviation is 3 minutes rather than 12, so its modeled times vary less around the same 30-minute mean. No A outcome exceeds 40 minutes, compared with $2/8=25\%$ for B. Equal means do not imply equal chances of being late, so the dispatcher's comparison hides the difference this customer cares about. This conclusion uses the eight equally likely outcomes in the stated model.}

\begin{summaryexitbox}{EXIT}
One thing the video changed about my first take: \hrulefill
\end{summaryexitbox}

\end{document}
